Foundation model-based multi-agent systems are moving into production, yet work on runtime reliability and serving efficiency is still fragmented across diagnostics and infrastructure. This review synthesizes 137 references, 45 systems, and 8 benchmarks through a unified runtime-stack perspective. We organize diagnostics into failure localization, observability, safety, and benchmarks, and infrastructure into memory management, prefix-aware scheduling, cache lifecycle, and power management. The central synthesis is that both layers act on overlapping runtime state but optimize different objectives. Diagnostics benefits from retaining traces, context, and replayable state, whereas serving infrastructure reclaims or compresses state to satisfy memory, latency, and throughput budgets. We term this interaction the \emph{diagnostics--infrastructure retention tension}. We further identify three workload characteristics---sparse activation, invocation distance, and inter-agent memory sharing---that shape agent-aware serving, and use explicit evidence tiers to separate full-text-verified results from abstract-level or contextual evidence. Rather than ranking heterogeneous systems, we derive testable co-design directions: retention-aware eviction, failure-aware scheduling, budgeted observability, and infrastructure-triggered diagnostics. This framing positions runtime state retention as a shared systems resource and clarifies where current multi-agent diagnostics and serving research remain disconnected.
